Tento dokument predstavuje detailnú dokumentáciu zameranú na návrh, konštrukciu,
realizáciu a oživenie hardvéru laboratórneho modelu pre inteligentnú reguláciu teploty
chladiacej kvapaliny. 
Model využíva uzavretý hydraulický okruh, v ktorom je hlavným
aktívnym prvkom Peltierov článok (Thermoelectric Cooler - TEC).

Cieľom hardvérového riešenia je zabezpečiť spoľahlivý zber termodynamických údajov
z expanznej nádržky a implementovať výkonové elektronické obvody pre precízne riadenie
akčných členov: vodná pumpa, Peltierov článok, ventilátor chladiča. 
Vzhľadom na zadanie implementácie troch pokročilých regulačných režimov (Mód 1: Regulácia cez TEC, Mód 2:
Regulácia cez prietok pumpy, Mód 3: Kaskádová/hybridná regulácia) bol kladený dôraz na
linearitu výkonových prvkov a bezpečnosť prevádzky. Kritickým aspektom celého návrhu
je manažment teplej strany Peltierovho článku, kde by bez adekvátneho odvodu tepla
došlo k okamžitému zniženiu komponentu.

Na záver by sme sa chceli úprimne poďakovať otcovi člena nášho tímu, Samuela Z., za jeho neoceniteľnú pomoc pri konštrukcii hardvérovej časti systému. 
Jeho odborné elektrotechnické znalosti a ochota zapožičať potrebné náradie výrazne prispeli k úspešnej fyzickej realizácii tohto projektu.